% 原文第39页；PCB铜箔图与器件位号保留。
\zhbox{54}{64}{235}{80}{9.5}{\textbf{PCB 布局指南}\\合理布局对稳定工作至关重要。不当布局会导致性能下降、EMI 增大、电阻损耗增加和系统不稳定。为获得最佳结果，请参阅图15并遵循以下指南：}
\zhbox{54}{146}{235}{57}{9.5}{1. 降压模式下，输入功率回路（$C_{IN}$、SWA、SWB）和逐周期电流检测电阻 $R_{SENSE1}$ 应尽量靠近放置。}
\zhbox{54}{206}{235}{76}{9.5}{2. 升压模式下，输出功率回路各元件，包括输出滤波电容 $C_{OUT}$、功率 MOSFET（SWC、SWD）和逐周期电流检测电阻 $R_{SENSE1}$，应尽量靠近放置。}
\zhbox{54}{287}{235}{26}{9.5}{3. 使用宽铜走线和功率回路过孔，以改善散热。}
\zhbox{54}{318}{235}{39}{9.5}{4. 将裸露焊盘连接至 GND，并在裸露焊盘上布置散热过孔。}
\zhbox{54}{362}{235}{26}{9.5}{5. 将小型去耦电容放在靠近 VIN、VOUT 和 AGND 的位置。}
\zhbox{54}{390}{235}{30}{9.5}{6. 栅极驱动走线及回流路径应尽量直接。}
\zhbox{54}{424}{235}{65}{9.5}{7. 正向和回流走线应彼此靠近，以尽量减小栅极驱动路径的电感。可将两条走线并排布置，或在相邻层上上下重叠布置。}
\zhbox{54}{495}{235}{76}{9.5}{8. 对平均电流检测电阻 $R_{SENSE2}$ 和逐周期电流检测电阻 R4 使用开尔文连接。从 $R_{SENSE2}/R_{SENSE1}$ 端子到 IC 引脚的这些线应平行走线。避免经过 SW1、SW2 或栅极驱动走线等噪声区域。}
\zhbox{54}{577}{235}{38}{9.5}{9. 将电流检测信号滤波电容尽量靠近 IC 引脚放置。}
\zhbox{54}{621}{235}{26}{9.5}{10. 将 VCC 和 AVDD 电容分别尽量靠近 VCC 和 AVDD 引脚放置。}
\zhbox{54}{652}{235}{39}{9.5}{11. 将 BST1 自举电容靠近 IC 放置，并直接连接至 BST1 和 SW1 引脚。}
\zhbox{324}{64}{235}{39}{9.5}{12. 将 BST2 自举电容靠近 IC 放置，并直接连接至 BST2 和 SW2 引脚。}
\zhbox{324}{108}{235}{51}{9.5}{13. 反馈环路走线应远离所有噪声源。将 FB 分压电阻 R1 和 R2 尽量靠近 FB 和 AGND 引脚放置。}
\zhbox{324}{166}{235}{36}{9.5}{14. 分开功率和信号路径，使功率电流或开关电流不会流经 AGND 连接。}
\zhbox{324}{209}{235}{51}{9.5}{15. 在 PGND 引脚附近、VCC 电容的 PGND 连接处附近，或 PGND 与 R4 的连接处附近，将 PGND 和 AGND 走线相连。}
\zhbox{520}{287}{31}{6}{4.5}{顶层}
\zhbox{520}{303}{31}{6}{4.5}{过孔}
\zhbox{520}{318}{31}{6}{4.5}{底层}
\zhbox{407}{474}{74}{15}{10}{\centering\textbf{顶层}}
\zhbox{405}{675}{76}{13}{10}{\centering\textbf{底层}}
\zhbox{348}{690}{190}{13}{9.5}{\centering\textbf{图15：推荐 PCB 布局}}

